\documentclass[11pt]{article}
\usepackage[margin=1.05in]{geometry}
\usepackage[T1]{fontenc}
\usepackage{lmodern,microtype,amsmath,amssymb,amsthm,mathtools,hyperref}
\hypersetup{colorlinks=true,linkcolor=blue,citecolor=blue,urlcolor=blue}
\numberwithin{equation}{section}
\newtheorem{theorem}{Theorem}[section]
\newtheorem{lemma}[theorem]{Lemma}
\newtheorem{proposition}[theorem]{Proposition}
\newtheorem{corollary}[theorem]{Corollary}
\theoremstyle{remark}\newtheorem{remark}[theorem]{Remark}
\newcommand{\E}{\mathbb E}
\newcommand{\Var}{\operatorname{Var}}
\newcommand{\Hyp}{\operatorname{Hyp}}
\newcommand{\artanh}{\operatorname{artanh}}
\newcommand{\argmax}{\operatorname*{arg\,max}}
\title{Balanced Populations Maximize the Optimal Sub-Gaussian Proxy\\for Hypergeometric Sampling at Every Finite Size}
\author{Zinan Wang}
\date{16 October 2026}
\begin{document}\maketitle
\begin{abstract}
For every population size $N\ge2$ and every fixed draw count $1\le m<N$,
we determine all marked-object counts $k$ maximizing the optimal centered
sub-Gaussian variance proxy of the hypergeometric distribution.
The complete set of maximizers is exactly
$\{\lfloor N/2\rfloor,\lceil N/2\rceil\}$.
Every noncentral success count lies strictly below the ordinary variance
at a central count, and an explicit uniform finite-population separation
is established. For even populations the maximal proxy is
$m(N-m)/[4(N-1)]$. For odd populations the maximizing count is known
exactly and its proxy exceeds the corresponding variance; an exact
one-dimensional finite-sum characterization is given, with elementary
closed forms for singleton and two-draw sampling.
The argument combines a sharp all-tilt curvature bound with a scaled
Bernoulli proxy inequality, in different sectors of the margin domain.
The proof is analytical, accompanied by exact-rational and numerical audits.
\end{abstract}
\noindent\textbf{Keywords.} Hypergeometric distribution; optimal sub-Gaussian
proxy; finite populations; exponential tilting; extremizer classification.

\section{Introduction}
Let $H_{N,k,m}$ count marked items in a size-$m$ sample without replacement
from an $N$-element population containing $k$ marked items. Its variance is
maximized at the central marked counts. It is not automatic that the same
holds for the optimal centered sub-Gaussian variance proxy: skewed bounded
distributions have optimal proxies strictly larger than their variances
\cite{Arbel}. We establish this centrality for all finite sample and population
sizes and give a strictly positive separation certificate for every off-center
margin. In particular, the previous fixed-$m$ eventual result in the
companion work \cite{Wang} is replaced by an all-$N$ theorem.

The proof imports two already proved facts from \cite{Wang}: the all-tilt
hypergeometric variance upper bound, and a scaled Bernoulli optimal-proxy
comparison. The present manuscript supplies the discrete global argument
linking these inputs, a quantitative margin gap, and its exact value consequences.
The dependence on those theorems is explicit. Classical concentration results
and optimal variance-proxy studies appear in \cite{Hoeffding,Serfling,GW,BM,Arbel,ThreeMass}.
No unpublished collaboration manuscript is used or compared.

\section{Distribution, proxy, and input bounds}
For $N\ge2$ and $1\le k,m<N$, define
\begin{equation}\label{pmf}
\Pr(H_{N,k,m}=j)=
 \frac{\binom{k}{j}\binom{N-k}{m-j}}{\binom{N}{m}},
 \quad \max(0,k+m-N)\le j\le\min(k,m).
\end{equation}
The mean is $\mu=km/N$, with
\begin{equation}\label{variance}
\Var(H_{N,k,m})=
M_{N,m}\frac{k(N-k)}{N^2},\qquad
M_{N,m}:=\frac{m(N-m)}{N-1}.
\end{equation}
Let
\begin{equation}\label{proxy}
\psi_{N,k,m}(t)=\log\E e^{t(H_{N,k,m}-km/N)},\qquad
S_{N,m}(k)=\sup_{t\in\mathbb R\setminus\{0\}}
 \frac{2\psi_{N,k,m}(t)}{t^2}.
\end{equation}
The quotient extends continuously at $t=0$ with value
$\Var(H_{N,k,m})$ and tends to zero at both infinite endpoints.
Define $c_N=\lfloor N/2\rfloor$ and
\begin{equation}\label{W}
W_{N,m}=M_{N,m}\frac{c_N(N-c_N)}{N^2}
=\begin{cases}M_{N,m}/4&N\text{ even},\\
 M_{N,m}(1-N^{-2})/4&N\text{ odd}.
\end{cases}
\end{equation}

\begin{proposition}[Proved analytic inputs from \cite{Wang}]\label{inputs}
The actual hypergeometric optimal proxy satisfies
\begin{equation}\label{C}
S_{N,m}(k)\le
\frac{\min\{k(N-k),m(N-m)\}}{4(N-1)}
\quad (1\le k,m<N).
\end{equation}
For $1\le m\le k<N/2$,
\begin{equation}\label{B}
S_{N,m}(k)\le M_{N,m}\,\beta(k/N),\quad
\beta(p)=\frac{1-2p}{2\log((1-p)/p)}.
\end{equation}
Equality holds in \eqref{B} if $m=1$; the underlying comparison
for $m\ge2$ is strict. Success complementation $k\mapsto N-k$
and sample complementation $m\mapsto N-m$ preserve $S$.
\end{proposition}
The upstream theorem names include
\texttt{hypergeomTiltVariance\_sharp} and the actual-proxy
curvature bridge, and the scaled Bernoulli bound in the
\texttt{OddHypergeomExtremizers} module. The unnormalized proxy $S$
is related to the companion normalized coefficient by multiplication
by $m(N-m)/(4N)$. Thus no transposition of a normalized expression
is required when invoking \eqref{C}--\eqref{B}.

\section{The exact finite extremizer theorem}
\begin{theorem}[All population sizes, all fixed sample sizes]\label{main}
For every integer $N\ge2$ and every $1\le m<N$,
\begin{equation}\label{argmax}
\argmax_{1\le k<N}S_{N,m}(k)
=\{\lfloor N/2\rfloor,\lceil N/2\rceil\}.
\end{equation}
Define the strictly positive quantity
\begin{equation}\label{delta}
\Delta_{N,m}=
\begin{cases}
\dfrac{M_{N,m}}{3N^2+4},& N\text{ even},\\[6pt]
\dfrac{M_{N,m}(2N^2-3)}{4N^2(N^2+3)},& N\text{ odd}.
\end{cases}
\end{equation}
Every noncentral count $k\notin\{c_N,N-c_N\}$ obeys
\begin{equation}\label{strict}
S_{N,m}(k)\le W_{N,m}-\Delta_{N,m}
<W_{N,m}\le S_{N,m}(c_N).
\end{equation}
The noncentral assertion is vacuous for $N=2,3$.
\end{theorem}

\begin{lemma}[Cubic lower bound for the inverse hyperbolic tangent]\label{atanh}
For $0<d<1$,
\begin{equation}\label{Bdeficit}
\artanh(d)\ge d+\frac{d^3}{3},\qquad
\frac{d}{4\artanh(d)}
\le\frac{1}{4(1+d^2/3)}.
\end{equation}
The first inequality is strict for $d>0$.
\end{lemma}
\begin{proof}
The difference on the left vanishes at zero and has derivative
$\frac{1}{1-d^2}-1-d^2=d^4/(1-d^2)>0$.
The second inequality follows by taking reciprocals of positive quantities.
\end{proof}

\begin{proof}[Proof of Theorem \ref{main}]
Complementation reduces the analysis to
$\bar k=\min(k,N-k)$ and $\bar m=\min(m,N-m)$, both in
$[1,\lfloor N/2\rfloor]$. Set $M=M_{N,m}$ and $W=W_{N,m}$,
and assume $\bar k<c_N$.

First suppose $\bar m\le\bar k$. Put $d=1-2\bar k/N>0$.
By \eqref{B} and Lemma \ref{atanh},
\begin{equation}\label{cap1}
S_{N,m}(k)\le\frac{M}{4(1+d^2/3)}.
\end{equation}
If $N$ is even, offcenter $\bar k\le N/2-1$ entails
$d\ge 2/N$, whence
\[
S_{N,m}(k)\le
\frac{3MN^2}{4(3N^2+4)}
=\frac{M}{4}-\frac{M}{3N^2+4}=W-\Delta_{N,m}.
\]
If $N$ is odd, $\bar k\le(N-3)/2$ entails $d\ge3/N$,
and
\[
S_{N,m}(k)\le\frac{MN^2}{4(N^2+3)}
=\frac{M(N^2-1)}{4N^2}
-\frac{M(2N^2-3)}{4N^2(N^2+3)}
=W-\Delta_{N,m}.
\]

Next suppose $\bar k<\bar m$. By the sharp curvature
ceiling \eqref{C},
\begin{equation}\label{cap2}
S_{N,m}(k)\le
\frac{\bar k(N-\bar k)}{4(N-1)}.
\end{equation}
The exact difference
\begin{equation}\label{disc}
\bar m(N-\bar m)-\bar k(N-\bar k)
=(\bar m-\bar k)(N-\bar m-\bar k)
\end{equation}
is at least one for even $N$ and at least two for odd $N$.
In the even case,
\[
W-S_{N,m}(k)\ge\frac{1}{4(N-1)}
>\frac{M}{3N^2+4}=\Delta_{N,m},
\]
using $M\le N^2/[4(N-1)]$.
In the odd case the central variance factor is $1-N^{-2}$, so
\[
W-S_{N,m}(k)\ge
\frac{2-\bar m(N-\bar m)/N^2}{4(N-1)}
\ge\frac{7}{16(N-1)}
>\frac{M(2N^2-3)}{4N^2(N^2+3)}.
\]
Indeed, the final quantity is smaller than $1/[8(N-1)]$.
Thus \eqref{strict} holds for every noncentral margin.

At $\bar k=c_N$, the zero-tilt limit in \eqref{proxy} is
$W_{N,m}$, proving $S_{N,m}(c_N)\ge W_{N,m}$.
Success complementation gives
$S_{N,m}(N-c_N)=S_{N,m}(c_N)$.
All remaining $k$ have a strict deficit, completing the
classification and the quantified bound.
\end{proof}

\begin{proposition}[Stronger individual-margin certificate]\label{individual}
Let $k$ be noncentral, $\bar k=\min(k,N-k)$ and
$\bar m=\min(m,N-m)$, and put $d=1-2\bar k/N$.
The explicit bound
\[
G_{N,m}(k)=
\begin{cases}
W_{N,m}-\dfrac{M_{N,m}}{4(1+d^2/3)},&\bar m\le\bar k,\\[8pt]
W_{N,m}-\dfrac{\bar k(N-\bar k)}{4(N-1)},&\bar k<\bar m
\end{cases}
\]
satisfies $G_{N,m}(k)\ge\Delta_{N,m}$ and
$S_{N,m}(k)\le W_{N,m}-G_{N,m}(k)$.
In the first branch $G_{N,m}(k)$ increases with
$|N-2k|$.
\end{proposition}
\begin{proof}
Each branch is exactly the upper proxy used in
\eqref{cap1} or \eqref{cap2}; the lower bounds were
verified in the two cases of Theorem \ref{main}.
\end{proof}

\section{Value at the maximizing margin}
\begin{proposition}[Even populations]\label{even}
For every even $N\ge2$, the maximum has the closed form
\[
\max_{1\le k<N} S_{N,m}(k)=\frac{m(N-m)}{4(N-1)}
\quad(1\le m<N).
\]
At $k=N/2$ the optimal proxy equals the ordinary variance.
\end{proposition}
\begin{proof}
The maximizing $k$ follows from Theorem \ref{main}.
The all-tilt ceiling \eqref{C} at this $k$ yields
$S\le M/4=W$; the zero-tilt limit yields $S\ge W$.
\end{proof}

\begin{proposition}[Odd populations]\label{odd}
For $N=2n+1\ge3$ and any $1\le m<N$,
$\max_{1\le k<N}S_{N,m}(k)=S_{N,m}(n)>W_{N,m}$.
Writing $\mu=mn/N$, this value is exactly the maximum
over all real tilts of the continuous-at-zero function
\[
F_{N,m}(t)=
\begin{cases}
\displaystyle\frac{2}{t^2}\log\!
\left(\sum_{j=\max(0,m-n-1)}^{\min(m,n)}
\frac{\binom{n}{j}\binom{n+1}{m-j}}{\binom{2n+1}{m}}
e^{t(j-\mu)}\right),&t\ne0,\\[10pt]
W_{N,m},&t=0.
\end{cases}
\]
Every maximizing nonzero tilt $t$ satisfies
$t\psi_{N,n,m}'(t)=2\psi_{N,n,m}(t)$.
No uniqueness or universal elementary closed form is asserted.
\end{proposition}
\begin{proof}
By complementation assume $m\le n$. The exact third
centered cumulant is
\[
\kappa_3=
W_{N,m}\frac{(N-2n)(N-2m)}{N(N-2)}
=W_{N,m}\frac{N-2m}{N(N-2)}>0.
\]
The Taylor formula
$2\psi(t)/t^2=W_{N,m}+\kappa_3t/3+O(t^2)$
is strictly above $W_{N,m}$ for small $t>0$.
The displayed finite sum follows from the mass function.
Its continuous extension tends to zero at both infinite ends,
so a finite maximizer exists; strictness ensures it is nonzero.
Differentiating there gives the stated necessary equation.
\end{proof}

\begin{corollary}[One- and two-draw closed forms]\label{small}
For odd $N\ge3$, the maximizing proxy for $m=1$ or $m=N-1$ is
\[
\frac1{2N\log((N+1)/(N-1))}.
\]
For odd $N\ge5$, the maximizing proxy for $m=2$ or $m=N-2$ is
\[
\frac2{N\log((N+1)/(N-3))}.
\]
\end{corollary}
\begin{proof}
For one draw use the optimal Bernoulli proxy
\cite{Arbel}. For two draws, specialize the exact three-mass
hypergeometric proxy in \cite{Wang,ThreeMass} and the
central success count established in Theorem \ref{main}.
Complementation supplies the two remaining cases.
\end{proof}

\section{Interpretation, computational audit, and formalization}
The theorem locates the maximum of the \emph{unnormalized}
optimal proxy for fixed $m$, not the separately normalized
coefficient in the global exchangeability optimization.
A central maximum set is also weaker than discrete unimodality:
no ordering of all consecutive success counts is claimed.
The explicit separation is not asserted to be best possible.
For odd $N$ and arbitrary $m\ge3$, Proposition \ref{odd}
is a finite one-dimensional optimization, not a closed
elementary evaluation or a theorem of uniqueness for the tilt.

The accompanying source code checks the exact rational
case inequalities independently of floating-point operations,
then directly optimizes the finite centered MGF with
log-sum-exp arithmetic. The frozen release passed 67,570
exact rational offcenter comparisons for $N\le60$, and
1,560 numerical proxy optimizations for $N\le18$.
The one- and two-draw closed forms were separately checked.
Numerical checks diagnose implementation errors; the
unbounded parameter conclusion rests solely on the proof.

The input results from \cite{Wang} and the present fixed-sample assembly
are kernel-checked in the companion Lean project. The public endpoints
\texttt{hypergeomOptimalProxy\_fixed\_sample\_extremizers},
\texttt{hypergeomOptimalProxy\_fixed\_sample\_gap\_even}, and
\texttt{hypergeomOptimalProxy\_fixed\_sample\_gap\_odd} cover the all-$N$ argmax and
explicit parity gaps. The odd central value is retained in its exact
finite-sum variational form; no elementary formula or uniqueness of the
optimizing tilt is used.

\section{Conclusion}
For every finite population and every fixed nontrivial sample size,
the closest-to-balanced success counts are exactly those
maximizing the optimal hypergeometric sub-Gaussian variance proxy.
An explicit positive gap excludes all other success counts
even by comparison with the central ordinary variance.
The classification is exact and uniform, whereas the general
odd-population proxy value remains a separate analytic
one-dimensional optimization.

\begin{thebibliography}{9}
\bibitem{Wang} Z.~Wang, \emph{Exact Hoeffding Constants for
Finite Exchangeable Sequences}, 2026,
\url{https://github.com/wzn1118/exact-hoeffding-constants}.
Precise proof inputs: branch \texttt{formal/full-closure-20261015},
commit \texttt{8687c3d}.
\bibitem{Arbel} J.~Arbel, O.~Marchal and H.~D.~Nguyen,
On strict sub-Gaussianity, optimal proxy variance and
symmetry for bounded random variables,
\emph{ESAIM: Probability and Statistics} \textbf{24}
(2020), 39--55.
\bibitem{ThreeMass} S.~Atouani, O.~Marchal and J.~Arbel,
Optimal sub-Gaussian variance proxy for 3-mass
distributions, arXiv:2510.06132 (2025).
\bibitem{Hoeffding} W.~Hoeffding, Probability inequalities
for sums of bounded random variables, \emph{JASA}
\textbf{58} (1963), 13--30.
\bibitem{Serfling} R.~Serfling, Probability inequalities
for the sum in sampling without replacement,
\emph{Ann. Statist.} \textbf{2} (1974), 39--48.
\bibitem{GW} E.~Greene and J.~Wellner, Exponential bounds
for the hypergeometric distribution, \emph{Bernoulli}
\textbf{23} (2017), 1911--1950.
\bibitem{BM} R.~Bardenet and O.-A.~Maillard,
Concentration inequalities for sampling without
replacement, \emph{Bernoulli} \textbf{21} (2015), 1361--1385.
\end{thebibliography}
\end{document}
